\section{Translation: electronic phenotyping as a coarsening problem}
\label{sec:translation}

\subsection{The data-generating process}

Consider a cohort of $n$ patients. Patient $i$ has a true clinical
state $\Y_i \in \{0,1\}$, where $\Y_i = 1$ indicates that the patient
has condition $Z$. The cohort prevalence is $\pi = \Prob(\Y_i = 1)$.

Each patient also has a latent severity $S_i$, a scalar summarizing how
advanced or symptomatic the condition is (for $\Y_i = 0$ patients, $S_i$
summarizes competing-condition burden or encounter intensity). Severity
is stochastically larger among true cases: $\E[S_i \mid \Y_i = 1] >
\E[S_i \mid \Y_i = 0]$.

The EHR records a diagnosis code for condition $Z$ through a coding
mechanism. Write $C_i \in \{0,1\}$ for the indicator that the code is
present. The coding mechanism is
\begin{equation}
  \Prob(C_i = 1 \mid \Y_i = 1, S_i)
    = \mathrm{logit}^{-1}(b_0 + b_1 S_i),
  \qquad
  \Prob(C_i = 1 \mid \Y_i = 0)
    = \mathrm{fpr},
  \label{eq:trans-coding}
\end{equation}
so that the code sensitivity for a true case rises with severity when
$b_1 > 0$, and code specificity is $1 - \mathrm{fpr}$. The marginal
code sensitivity is $\mathrm{sens} = \E[\,\mathrm{logit}^{-1}(b_0 + b_1
S) \mid \Y = 1\,]$. Additional coarse features, a procedure code or a
medication order, each carry their own sensitivity and specificity and
enter the same way.

The slope $b_1$ is the crux. When $b_1 = 0$, coding depends on $\Y$
only, not on severity, and the coarsening is non-informative: C2 holds.
When $b_1 > 0$, sicker true cases are preferentially coded; the
candidate set (the codes) carries information about the latent severity
beyond the latent state, and C2 fails. This is the
\emph{informative-coding} regime, and it is the realistic one: a true
case who is asymptomatic and never presents for care is rarely coded,
while a true case who is hospitalized is coded repeatedly.

The phenotyping problem is: given the observed codes $\{C_i\}$ for the
full cohort, and for a chart-reviewed subsample the adjudicated true
states $\{\Y_i\}$, recover the cohort prevalence $\pi$ and the coding
model.

\subsection{The translation}

\Cref{tab:translation} states the explicit correspondence to
masked-data inference.

\begin{table}[ht]
\centering
\small
\begin{tabular}{@{}p{0.40\linewidth} p{0.56\linewidth}@{}}
\toprule
\textbf{Masked-data framework} & \textbf{Electronic phenotyping} \\
\midrule
Component / latent cause
  & True clinical state / phenotype $\Y$ \\
Candidate set $c_i$
  & States consistent with the observed diagnosis codes \\
Masking probability
  & Coding mechanism $\Prob(\text{code} \mid \text{true state})$,
    shaped by clinician behavior and billing \\
Singleton candidate set $|c_i| = 1$
  & A chart-reviewed / clinician-adjudicated patient \\
C1 (truth in candidate set)
  & The true state is consistent with the codes;
    violated by gross miscoding \\
\textbf{C2} (coarsening at random): conditional on the candidate set, the masking distribution does not depend on which element is the true cause
  & Concretely: coding depends on $\Y$ only through modeled structure,
    not on prevalence or severity; \emph{violated} when coding
    intensity tracks severity (surveillance / detection bias) \\
C3 (parameter independence)
  & Coding-mechanism parameters distinct from prevalence;
    \emph{violated} when billing intensity tracks prevalence (a clinic
    facing a documented outbreak codes more aggressively, coupling
    sensitivity to the prevalence it is meant to estimate) \\
\bottomrule
\end{tabular}
\caption{Translation between the masked-data framework
\citep{towell2026masked} and electronic phenotyping. The C2 row is
load-bearing: informative coding is precisely the C2 violation.}
\label{tab:translation}
\end{table}

The phenotyping case is the binary specialization of the masked-cause
problem: the latent cause $\Y$ takes two values, and a patient with no
code at all has candidate set $\{0,1\}$ (the truth is unresolved),
while a chart-reviewed patient has a singleton candidate set
($\{0\}$ or $\{1\}$). The conceptual content is in the masking
mechanism. In the transcriptomics applications
\citep{towell2026scrnacoarsening,towell2026spatialcoarsening} the
coarsening mechanism is approximately non-informative once technical
covariates are modeled, so C2 holds and the difficulty is a rank
condition. In phenotyping, C2 fails by construction whenever $b_1 \neq
0$, and the difficulty is the bias that informative coding induces.

\subsection{Existing phenotyping methods in this language}

\paragraph{Rule-based phenotype algorithms} (the PheKB catalog,
\citealp{kirby2016phekb}) define a phenotype as a deterministic
function of codes, medications, and labs. In our language, a rule
declares a candidate set to be a singleton whenever the rule fires.
The rule's positive predictive value is exactly the probability that
this declared singleton contains the truth, a C1 statement.

\paragraph{PheWAS} \citep{denny2010phewas} treats the presence of a
code grouping as the phenotype directly. This is the code-only
estimator: it equates the candidate set with the truth and inherits
the full informative-coding bias.

\paragraph{Anchor-and-learn} \citep{halpern2016anchor} identifies
high-precision ``anchor'' observations whose presence almost certainly
implies the true state. An anchor is a near-singleton candidate set;
the framework predicts that anchors restore identifiability for the
same reason chart review does, with residual bias controlled by how
anchor-positive patients differ from the cohort.

\paragraph{Latent-class phenotyping} in the tradition of
\citet{hui1980estimating} and \citet{dawid1979maximum} fits
sensitivity and specificity jointly with prevalence, treating $\Y$ as a
latent class. This is the framework's natural estimator when no
singletons are available; \cref{thm:glass-ceiling} states the
conditions under which it is identifiable.

These methods are not in conflict. They are different ways of
supplying, or failing to supply, the identifying information that the
coarsening framework shows is required.
